% ===========================================================================
% Supplementary Appendix
% GRPO for Agentic LLM Fine-Tuning
% ===========================================================================

\documentclass{article}
\usepackage[utf8]{inputenc}
\usepackage[T1]{fontenc}
\usepackage{times}
\usepackage{url}
\usepackage{latexsym}
\usepackage{graphicx}
\usepackage{amsmath,amssymb}
\usepackage{booktabs}
\usepackage{microtype}
\usepackage{hyperref}
\usepackage{caption}

% ===========================================================================
% Supplementary Material
% ===========================================================================
\begin{document}

% ===========================================================================
\section{Complete Training Run Registry}
% ===========================================================================

\subsection{Tool-Use Experiments (Tinker Cloud GPU, Qwen3-8B)}

\begin{table}[htbp]
    \centering
    \caption{Complete tool-use GRPO runs on Tinker (Qwen3-8B, LoRA rank 32). All runs use Adam optimizer ($\beta_1$=0.9, $\beta_2$=0.95, $\epsilon$=1e-8), max sequence 1024 input / 512 generation.}
    \begin{tabular}{llcccccl}
        \toprule
        Run ID & Config & Group & Temp & Steps & Last-10 Reward & Status \\
        \midrule
        88ed2271 & Baseline: LR=3e-5 & 8 & 0.8 & 30 & 0.875 & Done \\
        2d488a85 & High LR: LR=1e-4 & 16 & 0.8 & 30 & 0.999 & Done \\
        5569c1fa & Low Temp: LR=3e-5 & 4 & 0.4 & 30 & 0.977 & Done \\
        22e9e5fd & xlam-60k: LR=3e-5 & 8 & 0.8 & 30 & 0.363 & Done \\
        386273f4 & Synthetic 5-tool & 4 & 0.8 & 100 & 0.825 & Done \\
        f63b618c & xlam-60k extended & 4 & 0.8 & 100 & 0.113 & Done \\
        1b01abef & MATH problems & 4 & 0.8 & 100 & 0.264 & Done \\
        \bottomrule
    \end{tabular}
    \label{tab:tool_runs}
\end{table}

\subsection{GSM8K Mathematical Reasoning Experiments (Tinker, Qwen3-8B)}

\begin{table}[htbp]
    \centering
    \caption{Complete GSM8K GRPO runs on Tinker (Qwen3-8B). Multi-seed replication and LoRA rank ablation series.}
    \begin{tabular}{llcccccc}
        \toprule
        Run ID & Seed & Rank & LR & Steps & Peak Acc & Last-10 Acc \\
        \midrule
        5db4e965 & 137 & 32 & 3e-5 & 50 & 62.5\% & 27.5\% \\
        aabb48cb & 256 & 32 & 3e-5 & 50 & 62.5\% & 32.5\% \\
        99971b26 & 512 & 32 & 3e-5 & 50 & 87.5\% & 30.0\% \\
        ba2a1694 & 42 & 8 & 3e-5 & 50 & 62.5\% & 21.2\% \\
        92ebcc48 & 42 & 16 & 3e-5 & 50 & 75.0\% & 18.8\% \\
        9219771c & 42 & 64 & 3e-5 & 50 & 87.5\% & 25.0\% \\
        1cc20cec & 42 & 32 & 5e-6 & 100 & 75.0\% & 27.5\% \\
        \bottomrule
    \end{tabular}
    \label{tab:gsm8k_runs}
\end{table}

\subsection{Google Colab QLoRA Experiments}

\begin{table}[htbp]
    \centering
    \caption{Local QLoRA experiments on Google Colab (T4 GPU, 16GB VRAM).}
    \begin{tabular}{lllccl}
        \toprule
        Model & Dataset & Method & LoRA Rank & Steps & Result \\
        \midrule
        Qwen2.5-0.5B & 14 synthetic & SFT only & 16 & 10 & JSON 30\% \\
        Qwen2.5-1.5B & Glaive-v2 & SFT+GRPO & 16 & 700 & JSON 92\% \\
        Qwen2.5-3B & ToolBench-v1 & SFT+GRPO & 8 & 40 & 0.91 multi-turn \\
        Llama-3.2-3B & GSM8K & GRPO only & 32 & 50 & 2.34\% acc \\
        Qwen3-4B & GSM8K & SFT+GRPO & 32 & 50 & In progress \\
        \bottomrule
    \end{tabular}
    \label{tab:colab_runs}
\end{table}

% ===========================================================================
\section{Training Hyperparameters}
% ===========================================================================

\subsection{Tinker Cloud GPU Configuration}

\begin{table}[htbp]
    \centering
    \caption{Tinker SDK hyperparameters for all GRPO runs.}
    \begin{tabular}{ll}
        \toprule
        Parameter & Value \\
        \midrule
        SDK Version & v0.16.1 \\
        Optimizer & Adam \\
        $\beta_1$ & 0.9 \\
        $\beta_2$ & 0.95 \\
        $\epsilon$ & 1e-8 \\
        Weight Decay & 0 \\
        Batch Size & 2 prompts/step \\
        Group Size & 4, 8, or 16 \\
        Max Input Length & 1024 tokens \\
        Max Generation Length & 512 tokens \\
        Warmup & None (constant LR) \\
        Checkpoint Frequency & Final only \\
        \bottomrule
    \end{tabular}
    \label{tab:tinker_hyperparams}
\end{table}

\subsection{LoRA Configuration}

\begin{table}[htbp]
    \centering
    \caption{LoRA adapter targets and ranks used across experiments.}
    \begin{tabular}{ll}
        \toprule
        Parameter & Value \\
        \midrule
        Target Modules & q\_proj, k\_proj, v\_proj, o\_proj, \\
                        & gate\_proj, up\_proj, down\_proj \\
        LoRA Ranks Tested & 8, 16, 32, 64 \\
        Default Rank & 32 \\
        Alpha & 2 $\times$ rank (default) \\
        Dropout & 0.0 \\
        QLoRA (Colab) & 4-bit NF4, bfloat16 compute \\
        \bottomrule
    \end{tabular}
    \label{tab:lora_config}
\end{table}

% ===========================================================================
\section{Failure Case Analysis}
% ===========================================================================

\subsection{Three-Billion Parameter Failure Mode}

The Llama-3.2-3B model on GSM8K (run not shown in main paper) achieved only 2.34\% accuracy after 50 GRPO steps, barely above the 1.56\% random baseline. Analysis of training logs reveals the root cause: \textbf{56\% of training steps had zero gradient} because all G=4 sampled completions received identical rewards (all correct or all incorrect).

This ``zero-loss'' phenomenon indicates the 3B model lacks the reasoning capacity to generate diverse completion quality within a group. When all completions are equally wrong, advantages are zero and no learning occurs. The 8B model (Qwen3-8B) shows only 16--28\% zero-loss rate, enabling effective gradient updates.

\begin{table}[htbp]
    \centering
    \caption{Zero-loss step percentage across model sizes (lower is better).}
    \begin{tabular}{lccc}
        \toprule
        Model & Parameters & Zero-Loss \% & Final Acc \\
        \midrule
        Llama-3.2-3B & 3B & 56\% & 2.34\% \\
        Qwen3-8B (seed 512) & 8B & 16\% & 30.0\% \\
        Qwen3-8B (seed 137) & 8B & 28\% & 27.5\% \\
        Qwen3-8B (seed 256) & 8B & 24\% & 32.5\% \\
        \bottomrule
    \end{tabular}
    \label{tab:zero_loss}
\end{table}

\subsection{Qwen3-30B-MoE Volatility Analysis}

The Qwen3-30B-MoE model reached 99\% peak GSM8K training-step accuracy in our internal runs but exhibited 2.43$\times$ higher step-to-step reward variance than the dense 8B model. Statistical testing (Levene's test) confirms this difference is significant: $F(1, 98) = 22.4$, $p = 7.0 \times 10^{-6}$.

This volatility stems from MoE routing: different experts activate for different tokens, creating non-smooth loss surfaces. The sparse activation pattern means small changes in the router can dramatically shift which experts process each token, amplifying training instability.

\begin{table}[htbp]
    \centering
    \caption{MoE vs. Dense training stability comparison.}
    \begin{tabular}{lccc}
        \toprule
        Model & Step Var (SD) & Peak Acc & Final Acc \\
        \midrule
        Qwen3-8B Dense & 0.12 & 87.5\% & 30.0\% \\
        Qwen3-30B-MoE & 0.29 & 99.0\% & 95.0\% \\
        Ratio (MoE/Dense) & 2.43$\times$ & --- & --- \\
        \bottomrule
    \end{tabular}
    \label{tab:moe_volatility}
\end{table}

\subsection{MBPP Evaluation Failure}

The MBPP (Mostly Basic Python Problems) benchmark produced 0\% pass rate across all runs due to an output parsing bug in the evaluation script. The script expected a specific JSON format for extracted answers but the model output format differed, preventing correct answer matching.

This represents a methodological failure: an invalidated benchmark should not be reported. Only HumanEval results (which used a working evaluator) are reported in the main paper. Future work should verify evaluator correctness before running experiments.

% ===========================================================================
\section{Additional Experimental Details}
% ===========================================================================

\subsection{Reward Function Weighting Rationale}

The tool calling reward (0.3 + 0.4 + 0.3) was designed to prioritize:
\begin{itemize}
    \item \textbf{Correct tool name (+0.4):} Selecting the right tool is the core judgment decision
    \item \textbf{Valid JSON (+0.3):} Required for downstream parsing
    \item \textbf{Arguments present (+0.3):} Functional tool use requires populated arguments
\end{itemize}

These weights were not validated via ablation. Future work should test alternative weightings, particularly increasing the tool-name weight to further emphasize selection judgment over format compliance.

\subsection{Evaluation Protocol Limitations}

As noted in the main paper, our primary evaluation metrics are computed on \textbf{training prompts with new rollouts}, not on held-out test sets. Specifically:
\begin{itemize}
    \item Tool-use rewards: Fresh sampling from the same prompt distribution
    \item GSM8K accuracy: Random subset of 8 training examples per step (batch=2, group=4)
    \item HumanEval: Standard pass@1 on the 50-example benchmark (test set)
\end{itemize}

The GSM8K evaluation is the most compromised: evaluating on training data with small group sizes (G=4) means the ``test'' examples overlap with training distribution. True held-out evaluation on the GSM8K test set (1319 examples) was not performed.

\subsection{Multi-Turn Tool Calling Reward Breakdown}

The multi-turn reward function (5-component) enables finer-grained learning signals:

\begin{table}[htbp]
    \centering
    \caption{Multi-turn tool calling reward components.}
    \begin{tabular}{ll}
        \toprule
        Component & Weight \\
        \midrule
        First turn = tool call & +0.25 \\
        Final natural language answer & +0.30 \\
        Arguments populated & +0.15 \\
        Clean JSON output & +0.10 \\
        Repeated tool call (penalty) & -0.30 \\
        \bottomrule
    \end{tabular}
    \label{tab:multiturn_reward}
\end{table}

The -0.30 penalty for repeated tool calls was critical: it eliminated SFT's looping failure mode entirely (Section 4.1.2 of main paper). Without this penalty, models would repeatedly call the same tool in an infinite loop.

Because these scores come from custom scenario rollouts rather than a standard benchmark, we emphasize that they are internal protocol numbers. We did not measure inter-rater reliability, did not use multiple human judges, and did not adopt proxy-state final-answer evaluation. Future revisions should benchmark the same agents under ToolRM / FC-RewardBench-style outcome evaluation or other standardized multi-turn frameworks.

% ===========================================================================
\section{Qualitative Examples}
% ===========================================================================

\subsection{Successful Tool Call (Post-GRPO)}
\begin{verbatim}
<|im_start|>user
What is the weather in Bangalore?<|im_end|>
<|im_start|>assistant
{
  "tool_calls": [
    {
      "name": "get_weather",
      "arguments": {
        "location": "Bangalore",
        "unit": "celsius"
      }
    }
  ]
}<|im_end|>
\end{verbatim}

\subsection{Failed Tool Call (Pre-GRPO, Plain Text)}
\begin{verbatim}
<|im_start|>user
What is the weather in Bangalore?<|im_end|>
<|im_start|>assistant
The weather in Bangalore is sunny and warm today.
The temperature is around 28 degrees Celsius.
It should be a nice day to go outside!<|im_end|>
\end{verbatim}

The pre-GRPO model generates natural language responses instead of structured tool calls. After GRPO training, the model reliably produces valid JSON with correct tool names and populated arguments.

\subsection{Two-Phase Learning Trajectory (Representative)}

\begin{verbatim}
Step 1:  Reward=0.00, Loss=0.85, Output="42" (format: plain text)
Step 5:  Reward=0.05, Loss=0.80, Output="The answer is 42." (format: text with number)
Step 10: Reward=0.10, Loss=0.72, Output="\boxed{42}" (format: boxed answer)
Step 15: Reward=0.14, Loss=0.65, Output="\boxed{42}" (stable format)
Step 20: Reward=0.40, Loss=0.55, Output="\boxed{42}" (reasoning begins)
Step 25: Reward=0.58, Loss=0.40, Output="\boxed{42}" (reasoning improves)
Step 30: Reward=0.55, Loss=0.38, Output="\boxed{42}" (plateau)
\end{verbatim}

This trajectory illustrates the two-phase pattern: format compliance (steps 1-20) precedes reasoning improvement (steps 20-30).

% ===========================================================================
\section{Platform and Compute Costs}
% ===========================================================================

\begin{table}[htbp]
    \centering
    \caption{Compute platform usage and estimated costs.}
    \begin{tabular}{llcl}
        \toprule
        Platform & GPU & Experiments & Est. Cost \\
        \midrule
        Tinker SDK v0.16.1 & Cloud A100 & 17 runs & \$30--40 \\
        Google Colab Pro & T4 16GB & 5 team members & \$10/person \\
        HuggingFace Hub & N/A & Model hosting & Free \\
        \bottomrule
    \end{tabular}
    \label{tab:compute}
\end{table}

Total estimated spend: $\sim$\$50/person, well within typical capstone project budgets.

We report GPU platform and wall-clock budget, but not a full token-processed accounting per run. Likewise, our data-split story is uneven across domains: HumanEval uses a held-out benchmark subset, while tool calling and GSM8K largely report training-distribution rollouts. This split mismatch is a major threat to cross-domain comparability.

% ===========================================================================
\section{Code and Reproducibility}
% ===========================================================================

All training scripts are available in the repository:
\begin{itemize}
    \item \texttt{grpo\_gsm8k\_base.py} -- Parameterized GSM8K training
    \item \texttt{grpo\_exp\_a\_baseline.py}, \texttt{exp\_b\_high\_lr.py}, \texttt{exp\_c\_low\_temp.py}, \texttt{exp\_d\_xlam.py} -- Tool-use experiments
    \item \texttt{grpo\_100\_xlam.py}, \texttt{grpo\_100\_synthetic.py}, \texttt{grpo\_100\_math.py} -- Extended runs
\end{itemize}

At present, we can release training scripts, paper sources, and the hardened held-out evaluation script, but the exact prompt packs / scenario definitions used for every internal tool-calling score are not yet packaged as a replication bundle. Releasing prompts, scoring rubrics, and scenario manifests is an explicit next step for reproducibility.

For stability diagnostics, we did not log KL-to-SFT trajectories or policy entropy, and our MoE analysis does not yet include routing entropy or expert-load imbalance measurements. Those omissions limit causal interpretation of both drift and volatility claims.

Training logs are available at: \texttt{/tmp/gsm8k\_*.log}, \texttt{/tmp/grpo\_*.log}

Model checkpoints are hosted on Tinker and HuggingFace Hub:
\begin{itemize}
    \item SWE-trained model: \texttt{huggingface.co/Madhu2133/qwen3-8b-swe-grpo}
    \item Tool-use checkpoints: \texttt{tinker://\textit{run\_id}/sampler\_weights/final}
\end{itemize}

\end{document}
